\section{Introduction}
\label{sec:introduction}

The planar circle-packing problem asks what fraction of the plane can be
occupied by congruent disks with disjoint interiors. Scale their radius
to $1/2$, so their centers form a set $\mathcal C\subset\RR^2$ with
$|x-y|\ge1$ for distinct centers. If $B_R$ is the disk of radius $R$
centered at the origin, we use the upper covered-area density
\[
 \overline\delta(\mathcal C)
 =\limsup_{R\to\infty}
 \frac{\vol\bigl(B_R\cap\bigcup_{c\in\mathcal C}B(c,1/2)\bigr)}
      {\vol B_R}.
\]
The triangular lattice attains $\pi/(2\sqrt3)$, and Thue's theorem gives
the matching upper bound. Hales presents an elementary geometric proof
based on an idea of Rogers~\cite[pp.~442--443]{Hales2000}.

The Fourier-analytic approach of Cohn and Elkies~\cite{CohnElkies2003}
asks whether the same optimum can be certified by one function with
compatible signs before and after Fourier transformation. The density
theorem alone does not supply such a function. Our result is a sharp
Fourier certificate for the classical optimum.

Throughout, the Fourier transform is
\[
  \widehat f(\xi)=\int_{\RR^2}f(x)e^{-2\pi i x\cdot\xi}\,\dd x.
\]
For real radial integrable functions it is real and radial. A Schwartz
function is a smooth function whose derivatives decay faster than every
inverse power of the radius; this regularity is stronger than the
admissibility conditions in the Cohn--Elkies theorem.

\begin{theorem}[Sharp Fourier certificate]\label{intro:main}
There is a real radial Schwartz function $f:\RR^2\to\RR$ such that
\begin{equation}\label{intro:conditions}
  \widehat f(0)=1,\qquad f(0)=\frac{2}{\sqrt3},\qquad
  \widehat f(\xi)\ge0\quad(\xi\in\RR^2),\qquad
  f(x)\le0\quad(|x|\ge1).
\end{equation}
\end{theorem}

Theorem~\ref{intro:main} gives an affirmative answer to the
two-dimensional case of Cohn and Elkies's Conjecture~7.3
\cite{CohnElkies2003}, concerning sharp auxiliary functions for their
linear programming bound. Their Definition~2.2 and Theorem~3.1 apply to
this Schwartz function and to arbitrary packings, not only periodic ones.
Their opposite Fourier phase gives the same transform for these even
functions.
For disks of radius $1/2$ they give
\[
  \vol B(0,1/2)\,\frac{f(0)}{\widehat f(0)}
  =\frac{\pi}{4}\frac{2}{\sqrt3}
  =\frac{\pi}{2\sqrt3}.
\]
This is the density of the lattice generated by $(1,0)$ and
$(1/2,\sqrt3/2)$. Thus the theorem identifies an extremizer for this
analytic bound. To put it in the equal-origin formulation of
Cohn and Elkies's Theorem~3.2, which their Conjecture~7.3 invokes, set
\[
 g(x)=\frac{\sqrt3}{2}\,
       f\bigl((\sqrt3/2)^{1/2}x\bigr).
\]
The two-dimensional scaling rule gives $g(0)=\widehat g(0)=1$,
$\widehat g\ge0$, and $g(x)\le0$ for
$|x|\ge(2/\sqrt3)^{1/2}$, exactly the requested planar witness.
The sign estimates also determine the zeros in their required domains.
For $f$, every zero with $|x|\ge1$ has integral squared radius. This
restriction recovers a classical equality consequence: among periodic
packings, equality in the density bound forces the triangular lattice
(Corollary~\ref{normalization:periodic-uniqueness}). The equal-origin
witness $g$ has zeros outside the shells of the triangular lattice of
covolume one, so it does not meet the additional zero-set requirement of
Cohn and Elkies's Conjecture~8.1. Section~\ref{sec:fourier} gives an explicit
extra node, and Section~\ref{sec:normalization} records the zero sets in
their respective sign domains and final radial coordinates.

\subsection{Fourier certificates and interpolation}

Cohn and Elkies formulated their sharpness conjecture in dimensions
$2$, $8$, and $24$~\cite[Conjecture~7.3]{CohnElkies2003}.
Their finite Laguerre--Gaussian approximations already combine prescribed
double roots, numerical searches for the remaining parameters, and rigorous
sign checks by rational arithmetic and Sturm root counts
\cite[Section~7, pp.~704--705]{CohnElkies2003}. Exact attainment of the
conjectured sharp bound remained unresolved in that approach.
Viazovska constructed the required function in dimension
$8$~\cite[Theorems~3 and~4]{Viazovska2017}, and Cohn, Kumar, Miller,
Radchenko, and Viazovska constructed its dimension-$24$ counterpart
\cite[Section~4]{CKMRV2017}.
Their constructions use modular forms to control a function and its
Fourier transform simultaneously.
Cohn and Elkies report that Gorbachev independently obtained essentially
the same linear programming bound \cite[p.~694]{CohnElkies2003}.

Interpolation gives a broader setting for these constructions.
Radchenko and Viazovska recovered even Schwartz functions on the real
line from their values and Fourier values at the square roots of
nonnegative integers~\cite[Theorem~1]{RadchenkoViazovska2019}.
Cohn, Kumar, Miller, Radchenko, and Viazovska developed interpolation
using values and first derivatives in dimensions $8$ and $24$ and used
it to prove universal optimality of the $E_8$ and Leech lattices
\cite[Theorems~1.4, 1.7, and~1.9]{CKMRV2022}.

The planar shell data have a different feature.
Cohn, Kumar, Miller, Radchenko, and Viazovska conjectured that values
and first derivatives of a radial Schwartz function and its Fourier
transform on the nonzero shells of the triangular lattice normalized to
covolume one do not determine
the function uniquely~\cite[Conjecture~7.5]{CKMRV2022}.
Sardari proved this nonuniqueness by constructing infinitely many
linearly independent radial Schwartz functions with double zeros,
before and after Fourier transformation, on those shells
\cite[Theorem~1.7]{Sardari2021}.
Thus prescribing the shell data alone does not single out a planar
certificate. We solve a specified interpolation problem within an
explicit family of Gaussian-damped cardinal functions and establish
the signs separately.

A 2026 bachelor's thesis abstract by Zhitniaia~\cite{Zhitniaia2026}
announces an integral representation of a magic function for the
hexagonal lattice using weakly holomorphic modular forms and theta
functions for $\Gamma_0(12)$. Our comparison is limited to the
institutional abstract, which does not specify the complete sign,
normalization, and regularity statement. The precise relation of that
claim to Theorem~\ref{intro:main} remains to be determined from the
full thesis.

\subsection{The construction}

For a sharp Fourier bound on a lattice, Poisson summation forces
the auxiliary function and its Fourier transform to vanish on the
nonzero shells of the lattice and its dual, respectively
\cite[Section~5]{CohnElkies2003}.
After normalizing the triangular lattice to covolume one, it is
isometric to its dual, so the two sets of shell radii coincide.
This normalization makes the coupled Fourier equations symmetric.
At the end we dilate the resulting function to put the first
physical-space shell at the exclusion radius one in
Theorem~\ref{intro:main}.
We cover their squared radii, in a fixed scaled coordinate, by a
periodic set of interpolation nodes.
Periodic supersets of the values of the triangular quadratic form
also occur in Sardari's construction
\cite[Section~6.1]{Sardari2021}; here a fixed periodic set permits
an explicit trigonometric cardinal family.

An interior zero of a differentiable function with a fixed sign must
also have zero first derivative. We therefore impose zeros of order
at least two on the Fourier side and at every physical-space shell
beyond the first. The first physical-space shell is the boundary of
the required sign region, so a simple zero is possible there; its
prescribed negative derivative will fix the normalization.

A trigonometric polynomial vanishes doubly on the chosen set.
Dividing it by linear and quadratic factors produces functions
that prescribe values and first derivatives at individual nodes.
This Hermite-cardinal pattern has a precedent in the Gaussian
Beurling--Selberg constructions of Carneiro, Littmann, and Vaaler
\cite[formulas~(1.2)--(1.3)]{CarneiroLittmannVaaler2013}; the Fourier
coupling and sign estimates below are proved for the present family.
Multiplication by a fixed Gaussian gives radial Schwartz functions.
Their Fourier transforms admit integral representations with
additional Gaussian decay. We use that decay to solve the coupled
interpolation equations: a finite matrix handles the first nodes,
and explicit estimates control the remaining infinite system. The inversion
acts on the full unweighted $\ell^1$ coefficient space. Its uniqueness
concerns the prescribed coefficient lists within this chosen family.

The resulting functions have exact zeros, but the coefficient
lists used for computation are finite approximations.
An explicit per-list $\ell^1$ error bound connects these approximations
to the exact solution. Near a zero, a raw error divided by its vanishing
power would be singular. We instead subtract the lower Taylor terms of
the error before division; positive Bernstein coefficients then certify
the signs on complete half-gaps. A separate sine-product barrier and
second-derivative estimate control the entire unbounded region.
Poisson summation first equates the two origin values.
Differentiating its dilation identity then determines their common
positive value from the prescribed first-shell derivative, so the
final normalization is exact.

Section~\ref{sec:fourier} develops the cardinal functions and their
Fourier transforms. Section~\ref{sec:interpolation} solves the
interpolation equations, and Section~\ref{sec:signs} proves the
global signs and identifies their zeros. Section~\ref{sec:normalization}
establishes the exact normalization, performs the final dilation, and
derives the packing consequences. Appendix~\ref{sec:certificates} gives
the finite data and the primary interval certificate, including the
analytic integration error that connects its finite sums to exact moments.
Appendix~\ref{ind:section} gathers the rational alternatives and
independent checks, with their precise mathematical and computational
scopes.
